\section{Unit-normalized signed moments and a single maximal depth}
\label{usm-section}

Let \(n>324\) be prime, \(B=n^4\), \(L=\log(8B)\), and
\[
 w_k=3k+\zeta,\qquad \alpha_k=w_k/\bar w_k,\qquad \zeta^2=\zeta-1,
 \qquad B\leq k<2B .
\]
Fix a subset \(I\) of this actual integer block on which the
\(\alpha_k\) are multiplicatively independent. Write
\[
 w_k^n=A_k+C_k\zeta,\quad T_k=|A_kC_k(A_k+C_k)|,\quad
 t_k=\log\max(|A_k|,|C_k|,|A_k+C_k|).
\]
The primitive and unramified power identities, together with the
individual-arm cap of Lemma~\ref{adp-arm-cap}, give
\begin{equation}\label{usm-height}
 N(w_k)<36B^2,\qquad
 n\log(3B)\leq t_k<nL,\qquad v_q(T_k)\log q\leq t_k .
\end{equation}
The last bound holds for primes \(q>8B\). In fact the three nonzero
arms are pairwise coprime, so every such prime divides exactly one
arm and its entire power divides that arm.

\begin{lemma}[A coherent unit choice at one actual prime]
\label{usm-unit}
Fix a prime \(q>8B\). For every \(k\in I\) with \(q\mid T_k\), there
is a global \(\tau_{q,k}\in\{1,\zeta^2,\zeta^4\}\), independent of
the precision, such that
\[
 \beta_{q,k}=\tau_{q,k}\alpha_k,\qquad
 \beta_{q,k}^n=1\pmod {q^e}\quad(1\leq e\leq v_q(T_k)).
\]
The \(\beta_{q,k}\) remain multiplicatively independent and are
ratios of unit-rotated actual roots to their conjugates. The
norm-one \(n\)-torsion group modulo \(q^e\) has at most \(n\)
elements in both the split and inert cases.
\end{lemma}
\begin{proof}
In the Eisenstein residue algebra the unique zero arm gives
\[
 \alpha_k^n=
 \begin{cases}
 1,&C_k=0,\\
 \zeta^2,&A_k=0,\\
 \zeta^4,&A_k+C_k=0.
 \end{cases}
\]
For the third identity use \(\zeta-1=\zeta^2\) and
\(\bar\zeta-1=-\zeta\). All denominators are units: a prime dividing
a primitive boundary arm does not divide the norm. The arm is
already determined modulo \(q\), and remains the same at every
available precision. Since \(\gcd(n,3)=1\), its global cube root
of unity has a unique inverse \(n\)-th root in this group. Choose
that root as \(\tau_{q,k}\).

The map \(v\mapsto v/\bar v\) on the six global units has image
exactly \(\{1,\zeta^2,\zeta^4\}\). Thus \(\beta_{q,k}\) has the
claimed actual representation. Cubing any integer-exponent
relation among the \(\beta_{q,k}\) eliminates all unit factors;
independence of the \(\alpha_k\) then recovers every exponent.

At precision one the norm-one group is cyclic: it is
\(\mathbb F_q^\times\) in the split case, and the norm-one subgroup
of \(\mathbb F_{q^2}^\times\) in the inert case. It has at most
\(n\) elements killed by \(n\). Since \(q\nmid n\), the \(n\)-th
power map is invertible on the principal \(q\)-kernels, and each
such element has exactly one lift through each successive
precision. This proves the same cardinal bound at every level.
\end{proof}
The choices in Lemma~\ref{usm-unit} are coherent for each fixed
\(q\), but can change when \(q\) changes. They do not define one
common normalized map for all primes.

\begin{theorem}[An actual signed-ball injection]\label{usm-ball}
Let \(q>8B\) be prime and let \(e,\nu\) be positive integers with
\(e\log q\geq2\nu L\). Put
\[
 I_e(q)=\{k\in I:q^e\mid T_k\},\qquad M_e=|I_e(q)|.
\]
Then
\begin{equation}\label{usm-ball-count}
 C(M_e,\nu):=\sum_{j=0}^{\min(M_e,\nu)}
    2^j\binom{M_e}{j}\binom{\nu}{j}\leq n .
\end{equation}
\end{theorem}
\begin{proof}
For an integer vector \(x=(x_k)\) with \(\sum_k|x_k|\leq\nu\),
represent \(\prod_k\beta_{q,k}^{x_k}\) by \(z_x/\bar z_x\), using
the unit-rotated \(w_k\) for positive exponents and its conjugate
for negative exponents, with their indicated multiplicities.
The empty product is \(z_0=1\). These are nonzero actual
Eisenstein integers with
\[
 |z_x|=\sqrt{N(z_x)}\leq(6B)^\nu .
\]
Units and conjugation preserve the norm. For \(z=R+S\zeta\) and
\(z'=R'+S'\zeta\), the exact integer Gram identity is
\begin{equation}\label{usm-gram}
 4N(z)N(z')-(2RR'+RS'+SR'+2SS')^2
       =3(RS'-R'S)^2 .
\end{equation}
Consequently the determinant \(D=RS'-R'S\) of two representatives
obeys
\begin{equation}\label{usm-determinant}
 |D|\leq\frac2{\sqrt3}|z_x||z_y|
 \leq\frac2{\sqrt3}(6B)^{2\nu}<(8B)^{2\nu}.
\end{equation}
Indeed \((2/\sqrt3)(3/4)^{2\nu}\leq9/(8\sqrt3)<1\) for every
\(\nu\geq1\). This includes unequal lengths and empty products.

If the two residue products agree modulo \(q^e\), their conjugate
cross difference is divisible by \(q^e\). The factor
\(\zeta-\bar\zeta\) is a unit at \(q>3\), so \(q^e\mid D\).
The strict bound \eqref{usm-determinant} and the precision
assumption force \(D=0\). The exact algebraic ratios are equal,
and Lemma~\ref{usm-unit} recovers \(x=y\). Thus the actual integer
\(\ell^1\) ball injects into the norm-one \(n\)-torsion group.

To count this ball, select its \(j\) nonzero coordinates and
their signs. The positive absolute values with sum at most
\(\nu\) have cardinality \(\binom{\nu}{j}\), by subtracting one
from each and adjoining a slack coordinate. The \(j=0\) term
counts the zero vector, including \(M_e=0\). This gives
\eqref{usm-ball-count}. The slack coordinate is only part of
the counting argument; it is not an extra independent root.
\end{proof}
Expanding this finite sum gives
\begin{equation}\label{usm-special-counts}
 \begin{gathered}
 C(M,2)=2M^2+2M+1,\qquad
 3C(M,3)=4M^3+6M^2+8M+3,\\
 C(2,\nu)=2\nu^2+2\nu+1 .
 \end{gathered}
\end{equation}
The formulas include the small-dimensional and zero-radius cases
where applicable.

\begin{theorem}[The complete single-prime budget]\label{usm-budget}
Put \(r=\lceil\sqrt{n/2}\rceil\). For every prime \(q>8B\),
\begin{equation}\label{usm-five}
 \begin{gathered}
 \sum_{k\in I}(v_q(T_k)-3)_+\log q\leq\frac52 nL,\\
 \frac1B\sum_{k\in I}
       \frac{(v_q(T_k)-3)_+\log q}{t_k}\leq\frac5B .
 \end{gathered}
\end{equation}
There is no upper restriction on \(q\).
\end{theorem}
\begin{proof}
One has \(2r^2\geq n\) and \(r\leq\sqrt n\): use
\(\sqrt{n/2}\leq(3/4)\sqrt n\) and \(1\leq\sqrt n/4\).
Write \(w=\log q\) and \(h=\max(4,\lceil2rL/w\rceil)\).
Theorem~\ref{usm-ball} at \((e,\nu)=(4,2)\) gives
\(M_4\leq\sqrt{n/2}\). At \((e,\nu)=(h,r)\), two indices would
give \(n\geq C(2,r)=2r^2+2r+1>n\). Thus \(M_h\leq1\).

All layers from four through \(h-1\) cost at most
\[
 (h-4)M_4w\leq2r\sqrt{n/2}L
       \leq\sqrt2\,nL<\frac32 nL .
\]
The subtraction of four absorbs the ceiling error; the range is
empty when \(h=4\). The remaining layers have at most one actual
index, whose entire depth costs at most \(nL\) by
\eqref{usm-height}. This proves the first inequality.
Since \(L\leq2\log(3B)\), the lower height in
\eqref{usm-height} proves the second.
\end{proof}

\begin{theorem}[One actual maximum and two support thresholds]
\label{usm-one-owner}
For a prime \(q>8B\), set
\[
 h_0=\max(4,\lceil6L/\log q\rceil),\qquad
 h_1=\max(h_0,\lceil2rL/\log q\rceil).
\]
Let \(O_q\) contain one index of maximal actual depth, or be empty
if \(I\) is empty. Then
\begin{equation}\label{usm-remainder}
 \begin{split}
 \sum_{k\in I\setminus O_q}(v_q(T_k)-3)_+\log q
 &\leq(6M_4+2rM_{h_0})L\\
 &\leq7n^{5/6}L .
 \end{split}
\end{equation}
Its normalized remainder is at most \(14n^{-1/6}/B\) per prime.
If \(\log q\geq(r/2)L\), at most one root of \(I\) has depth at
least four. If \(\log q\geq2rL\), at most one root of \(I\) has
any positive depth.
\end{theorem}
\begin{proof}
The signed counts give
\[
 M_4\leq\sqrt{n/2},\qquad M_{h_0}<n^{1/3},\qquad M_{h_1}\leq1.
\]
For the middle inequality use
\((4/3)M^3<C(M,3)\leq n\). Maximality implies that \(O_q\)
contains every \(h_1\)-deep index. For each remaining depth
\(e<h_1\), the complete truncated inequality is
\[
 (e-3)_+\leq(h_0-4)\boldsymbol1_{e\geq4}
              +(h_1-h_0)\boldsymbol1_{e\geq h_0}.
\]
The two ceiling inequalities
\((h_0-4)\log q\leq6L\) and
\((h_1-h_0)\log q\leq2rL\) prove the first bound in
\eqref{usm-remainder}, including empty ranges. For the second,
\[
 6M_4+2rM_{h_0}\leq5\sqrt n+2n^{5/6}\leq7n^{5/6}.
\]
The height lower bound gives the normalized statement. Adding
back the entire selected index costs at most \(nL\).

At \(\log q\geq(r/2)L\), both thresholds equal four, hence the
positive-excess support has at most one index. At the stronger
threshold, Theorem~\ref{usm-ball} applies with \(e=1,\nu=r\),
which bounds the full positive-depth support by one.
\end{proof}
The first threshold allows other roots of depth one, two or three;
the second excludes them. Neither bounds the number of different
prime labels that can occur on one root.

\paragraph{Prime summation and the remaining problem.}
Every supported \(q>8B\) has rank \(n\), hence \(q\equiv1\) or
\(-1\pmod {3n}\). The same primary Brun--Titchmarsh bound used in
the adaptive-precision section, together with \eqref{usm-five},
gives, for real \(Z>8B\),
\begin{equation}\label{usm-progression}
 \frac1B\sum_{k\in I}
 \frac{\sum_{8B<q\leq Z}(v_q(T_k)-3)_+\log q}{t_k}
 \leq\frac{10Z}{B(n-1)\log(Z/(3n))}
 \leq\frac{11Z}{Bn\log(Z/(3n))}.
\end{equation}
For the actual finite set
\[
 \mathcal D_Z(I)=\{q>Z:q\text{ prime},\
                  \exists k\in I,\ v_q(T_k)\geq4\},
\]
the far positive mean is at most
\(5|\mathcal D_Z(I)|/B\). The needed bound
\(|\mathcal D_Z(I)|=o(B)\) has not been proved. The constants and
support thresholds improve; this result alone does not enlarge
the known useful asymptotic prime window.

In the strongest support region, a prime cannot supply both a
high-depth debit at one root of \(I\) and a depth-one or depth-two
credit at another root of \(I\). Credits at different primes and
outside \(I\) remain possible. The unit choices also depend on the
prime. Repeating the fixed-prime argument therefore supplies
neither one common finite target nor a bound for the full signed
far tail. Private-label counts and weights, pointwise exceptions,
and the independent-domain complement remain separate problems.

\paragraph{The finite arithmetic formally checked here.}
The module
\texttt{SignedMoment\allowbreak Arithmetic.lean} contains
22 new theorems, freshly compiled with its 27 unchanged local
dependencies under the pinned Lean \(4.32.0\) and Mathlib snapshot.
Every declaration has a complete axiom query using only the
standard three axioms. Its \texttt{ballCount} is the actual finite
binomial sum: its special formulas, monotonicity and numerical
consequences are proved. Interpreting that sum as an arithmetic
torsion image remains an explicit input to this module.

The actual finite depth filters, a singleton of maximal depth,
natural truncated subtraction, every remaining positive layer and
the per-index whole-depth caps enter the final finite theorems.
They prove the \(5/2\) total budget, the exact two-range remainder,
and the normalized \(5/B\) bound using individual heights
\(t_i\geq t_0>0\) with \(nL\leq2t_0\). A separate numerical bridge
accepts the concrete two-coordinate diamond lower bound.
For a completely natural threshold the module uses
\(r=\lfloor\sqrt{\lfloor n/2\rfloor}\rfloor+1\); it proves
\(2r^2\geq n\) and \(r^2\leq n\) for \(n\geq18\).
This is not asserted to equal the real ceiling in every even
perfect-square case. The polynomial inequality
\[
 4rM_4\leq r^2+4M_4^2\leq3n
\]
proves the finite \(5/2\) budget without a real square-root input.
The coefficient change from \(10/(n-1)\) to \(11/n\) is also
formally checked. The \(7n^{5/6}\) simplification, prime
distribution, actual residue torsion, and global tail membership
are not claimed as formal consequences of this numerical module.
